\begin{frame}
    \frametitle{Ejercicio}

    Determina si los vectores
    \[
        v_1=(1,0,1),\qquad
        v_2=(0,1,1),\qquad
        v_3=(1,1,1)
    \]
    forman una base de \(\mathbb{R}^3\).

    \medskip

    Resuelve el problema de dos formas:

    \begin{enumerate}
        \item mediante eliminación gaussiana;
        \item calculando un determinante.
    \end{enumerate}
\end{frame}


\begin{frame}
    \frametitle{Ejercicio}

    Determina para qué valores de \(a\in\mathbb{R}\) los vectores
    \[
        v_1=(1,0,1),\qquad
        v_2=(0,1,1),\qquad
        v_3=(1,1,a)
    \]
    forman una base de \(\mathbb{R}^3\).

    \medskip

    Para los restantes valores de \(a\), determina:

    \begin{enumerate}
        \item una relación de dependencia lineal;
        \item cuál de los vectores puede eliminarse sin cambiar
        las combinaciones lineales que se pueden obtener.
    \end{enumerate}
\end{frame}
